\section{Dataset and Evaluation Method}
\label{sec:method}

\subsection{Images and generated copies}

The paper's evaluation data is not public, so we built our own. The base corpus is 3,000 real
NFT images from a public Kaggle collection~\cite{kaggle_nft}: 1,000 each from Azuki, Bored Ape
Yacht Club (BAYC) and CryptoPunks. Azuki and BAYC are ordinary artwork, while CryptoPunks are
$24\times24$ pixel art. A separate 10,000-image CryptoPunks set~\cite{kaggle_punks} was kept
aside for stress tests and not mixed in, so that one collection does not dominate.

The images carry no edit labels, so we generated copies ourselves, following the six edit
types the paper takes from OpenSea's copymint examples (Figure~\ref{fig:edits}): flip,
rotate or mirror; crop, resize or reposition; a text, logo or emoji overlay; a
background-colour change; pixelation; and a colour swap or saturation change. We also generated exact copies as a control and
non-matching pairs as negatives. Each edit is applied with three random variations per base
image, using PIL.

\paragraph{One labelling caveat.} Our generator's \texttt{flip\_rotate\_mirror} edit applies an
optional mirror and then \emph{always} a rotation of 1--359 degrees. The category therefore
contains no pure mirrors, and every result for it describes rotation. We separate the two
in Section~\ref{sec:results}.

\begin{figure}[ht]
\centering
\begin{tikzpicture}[font=\footnotesize]
  \foreach \lab [count=\i from 0] in {original, exact copy, flip / rotate, crop / reposition,
                                      text / logo, background, pixelated, recoloured} {
    \pgfmathsetmacro{\x}{mod(\i,4)*3.6}
    \pgfmathsetmacro{\y}{-floor(\i/4)*1.15}
    \node[draw=ink2, rounded corners=2pt, minimum width=33mm, minimum height=8mm,
          fill=rule!40] at (\x,\y) {\lab};
  }
\end{tikzpicture}
\caption{The edit categories in our generated dataset: the six edit types from the paper, plus
exact copies as a control. Non-matching pairs serve as negatives.}
\label{fig:edits}
\end{figure}

\subsection{Train/test split}

The split is made \textbf{by base image, before any edit is applied}, so no image's copies
appear in both training and testing. We use 80\% for training and 20\% for testing: 2,400
training bases (48,000 generated images, 55,200 labelled pairs) and 600 test bases (12,000
images, 13,800 pairs). The paper uses the opposite split (20/80), because it tunes only a few
scalar thresholds. We train a Random Forest over 93 features and eight labels, which benefits
from more training data. The test split contains \textbf{11,400 copies and 2,400 non-copies}.

\subsection{The authors' reference set}

The paper's authors shared with us a small set of their own manipulated images: 202
manipulated CryptoPunks and 200 originals, with a 1,802-row file of labelled
original--candidate pairs (202 copies and 1,600 manipulated non-copies) and their own
precomputed hashes and distances. We use it for two things: to check that our reimplementation
matches theirs, and as an independent test made by a different generator. The set was shared
privately, so it stays local and is not redistributed.

\subsection{Pairwise evaluation}

Each test row is a pair (stored original, query) with a label saying whether the query is a
copy of that original. Every signal is scored on the pair directly, with no gallery or index in
the accuracy path. This matches how the dataset is labelled and how the hashes are evaluated.
It also rules out the failure described in Section~\ref{sec:oldbaseline}, where an index
decided the result instead of the signal.

\subsection{Metrics and operating points}

We report precision $P=\frac{TP}{TP+FP}$, recall $R=\frac{TP}{TP+FN}$ and
$F_1=\frac{2PR}{P+R}$. Two rules keep these honest:
\begin{itemize}
  \item \textbf{A recall always comes with its precision.} The single-edit test set is 82.6\%
        copies, so ``flag everything'' scores F1 90.5. On the geometric subset used for the
        sHash comparison (60\% copies) it scores 75.0. Detectors are compared at equal
        precision wherever possible.
  \item \textbf{Thresholds are chosen on training data and frozen for test.} The one exception
        is deliberate and works against us: in the head-to-head comparison sHash's threshold
        is chosen on the test set itself.
\end{itemize}
Inside the full detector, every signal is placed at the same false-alarm budget, at most 10\%
of training non-copies flagged: aHash $\le 7$, pHash $\le 19$, hsvHash $\le 3$, sHash $\le 21$,
ORB $> 23$ inliers. ORB's $> 16$ cutoff from the signal study sits at about 14\% false alarms,
so using it inside the detector would give ORB a looser vote than the others. We also evaluate
the detector at the paper's own thresholds (Section~\ref{sec:results}).

\subsection{Retiring an earlier ORB baseline}
\label{sec:oldbaseline}

The project inherited a first ORB pipeline that reported F1 73.6\% (precision 82.7\%, recall
66.3\%). Before building on it we checked it, and found that its gallery held 2,000 of the
3,000 originals: every CryptoPunk was missing. At $24\times24$ pixels ORB found no keypoints
in any punk, and the indexing code skipped images without descriptors without a warning. The
result was decided by which collection an image belonged to, not by whether it was a copy. A
script that does no image matching at all, only checking whether a filename is in the
gallery, reproduces all four cells of the reported confusion matrix exactly (30,210 TP,
15,390 FN, 6,306 FP, 3,294 TN). The threshold sweep was also flat (F1 73.58 to 73.52 across
every setting), which confirmed that the ORB score was not deciding anything. We retired those
numbers and rebuilt the evaluation pairwise, with the resolution normalisation of
Section~\ref{sec:architecture}.

\subsection{Reproducing the original hashes}
\label{sec:reimpl}

Before changing the detector, we reimplemented the four hashes in C11, plus dHash which sHash
uses internally. All five match Buchner's \texttt{imagehash} exactly on 600 distinct images,
including all 402 of the authors' reference images, high-resolution NFTs, native CryptoPunks
and synthetic edge cases. For sHash we matched the full ordered list of segment hashes, not
only the final distance. Our sHash distance also reproduces all 1,802 distances in the
authors' reference file. Two details made exact parity hard. First, Pillow declares some
variables as \texttt{float} but computes with \texttt{double} literals, so a natural all-float
port differs on about 0.4\% of pixels. Second, the paper's description of hsvHash differs from
the code that produced its numbers, as noted in Section~2.1.
